% ============================================================================
% CHAPTER 5: CONCLUSIONS
% ============================================================================
\chapter{Conclusions and Future Work}
\label{ch:conclusions}

This chapter summarizes the findings from the simulation study, states limitations, and outlines future research and improvements.

\section{Summary of Findings}
\label{sec:summary}

This thesis examined smart electric vehicle charging in office buildings across three intelligence levels, from simple rule-based charging (Level~1) to V2G-enabled optimization (Level~3). The following subsections condense the main results.

\subsection{Intelligence Level Comparison}
\label{subsec:intelligence_comparison}

\subsubsection{Level 1: Simple Rule-Based}

Level~1 reaches a daily cost of €31.63, which serves as the baseline. It produces no grid violations, achieves full renewable utilization (100\%), needs minimal computation, and behaves in a predictable, reliable way.

\textit{Verdict:} Suitable for small-to-medium installations (5--15~EVs) when simplicity and reliability matter most.

\subsubsection{Level 2: Intelligent Optimization}

Level~2 lowers daily cost to €30.45, a 3.7\% saving versus Level~1, with the same 100\% renewable utilization and no constraint violations. It relies on RL (Q-Learning) and therefore requires a training phase.

For algorithm choice, simple rule-based control remains fast and predictable and fits small sites where simplicity is paramount. Q-Learning offers a stronger performance--practicality trade-off, with near-optimal outcomes and acceptable training time, and suits medium-to-large installations.

\textit{Verdict:} Q-Learning is the primary recommendation where the fleet exceeds about ten vehicles; the simple rule-based controller remains adequate for smaller fleets.

\subsubsection{Level 3: Vehicle-to-Grid}

Level~3 shows €0.00 daily benefit and 0\% improvement versus simple charge-only, with €0.00 additional V2G revenue per day. Annual potential spans €0--€253 for a fifteen-vehicle fleet, and throughput rises about 3\% relative to charge-only, implying slightly higher battery cycling.

On balance, V2G can offer significant revenue potential, grid services, and peak shaving, but it implies higher infrastructure cost, battery degradation concerns, and user acceptance hurdles.

\textit{Verdict:} V2G becomes economically attractive when sell-back prices exceed roughly 80\% of purchase prices, price volatility is high enough for arbitrage, users tolerate extra cycling (e.g.\ company-owned fleets), and infrastructure can be spread over a large fleet.

\section{Answering the Research Questions}
\label{sec:research_questions}

\subsection{Which Intelligence Level Provides Substantial Benefits?}

\textit{Answer:} Level~1 (simple rule-based control) provides a basic implementation that is easy to maintain.

It finishes within 3.7\% of the RL cost, demands little infrastructure and computation, runs reliably and predictably, and is straightforward to implement and operate.

Level~2 adds incremental savings (3.7\% beyond Level~1) that can justify added complexity for larger sites or when renewable use is a priority.

\subsection{Is V2G Worth the Investment?}

\textit{Answer:} The case is context-dependent, though the setting remains promising for office buildings.

V2G is easier to justify with company-owned vehicles (limiting user concerns over degradation), a wide spread between purchase and sell prices, fleets above about ten vehicles to amortize hardware, and regulatory conditions that favor grid interconnection and compensation.

\section{Challenges and Limitations}
\label{sec:challenges}

\subsection{Technical Challenges}
\label{subsec:technical_challenges}

\subsubsection{Forecasting Accuracy}

Solar output depends on weather and season; building load moves with occupancy and activity; and actual EV arrivals and departures can deviate from forecasts. Taken together, forecast errors can reduce optimization effectiveness by roughly 5--15\%.

\textit{Mitigation:} Use robust optimization, model forecast uncertainty, and allow real-time updates.

\subsubsection{Computational Complexity}

Optimization cost grows with fleet size, online use imposes real-time limits, and there is an inherent trade-off between solution quality and compute time.

\textit{Mitigation:} Rely on approximate methods such as RL, parallelize computation, or push work to edge nodes.

\subsubsection{Communication and Control}

Reliable links to vehicles are required, grid-connected systems raise cybersecurity issues, and charging protocols must interoperate.

\textit{Mitigation:} Adopt established standards (OCPP, ISO~15118) and secure communication practices.

\subsection{Economic Challenges}
\label{subsec:economic_challenges}

\subsubsection{Electricity Market Structure}

Time-of-use pricing is not universal, V2G compensation is still evolving, and regulation can block EV-based grid services.

\textit{Mitigation:} Engage in policy discussion, join pilots, and explore alternative revenue streams.

\subsection{User Acceptance Challenges}
\label{subsec:user_challenges}

\subsubsection{Battery Degradation Concerns}

Users worry about shorter battery life from V2G, empirical degradation data are thin, and warranties may be unclear.

\textit{Mitigation:} Offer degradation safeguards or insurance, cap discharge aggressively, and support decisions with measured data.

\subsection{Study Limitations}
\label{subsec:study_limitations}

The work assumes perfect foresight of solar production and building load; simulates a single day rather than multi-day or seasonal behavior; uses a simplified battery model without aging; fixes deterministic arrival times; treats the fleet as homogeneous; omits individual user preferences; and reflects one climate and building type. Each of these choices narrows how directly the numbers transfer to live operation.

\section{Future Research Directions}
\label{sec:future_work}

\subsection{Short-Term Research Priorities}
\label{subsec:short_term}

\subsubsection{Uncertainty Quantification}

Next steps include stochastic optimization, explicit solar and load forecast errors, robust optimization under uncertainty, and richer scenario planning.

\textit{Expected impact:} More realistic performance estimates and policies that remain effective under uncertainty.

\subsubsection{Battery Degradation Modeling}

Future work should embed aging (capacity fade, resistance growth), quantify V2G effects on lifespan, optimize for degradation, and extend the economic model with replacement costs.

\textit{Expected impact:} Sharper V2G business cases and degradation-aware control.

\subsection{Medium-Term Research Directions}
\label{subsec:medium_term}

\subsubsection{Multi-Day and Seasonal Analysis}

Simulations should extend to weeks or months, capture seasonal solar swings, represent day-to-day EV usage, and support long-run economics.

\subsubsection{Heterogeneous Fleet Management}

Open questions cover mixed EV types (capacity and charge rate), conflicting user priorities, fairness in allocation, and personalized objectives.

\subsubsection{Advanced Forecasting}

Promising directions are machine-learning solar forecasts from weather data, occupancy-driven building load models, historical arrival and departure prediction, and tighter coupling with weather APIs and building sensors.

\subsection{Long-Term Research Vision}
\label{subsec:long_term}

\subsubsection{Advanced AI Techniques}

Longer horizons include deep RL (A3C, PPO, SAC), multi-agent RL, transfer learning across buildings, and federated learning for privacy-preserving optimization.

\subsection{Emerging Technologies}
\label{subsec:emerging_tech}

\subsection{Environmental Impact}
\label{subsec:environmental}

Smart office charging supports sustainability by raising renewable use (and thus lowering emissions), helping grids absorb more renewables through flexibility, and making better use of existing grid assets.

\subsection{Social Impact}
\label{subsec:social_impact}

Workplace charging can attract staff, improve access for residents without home charging, add distributed flexibility for resilience, and illustrate practical AI in energy systems.

\section{Final Remarks}
\label{sec:final_remarks}

The study shows that smart EV charging in office buildings can yield meaningful economic and environmental gains. RL improves costs by a further 3.7\% over the simple baseline. Returns diminish by level: Level~2 delivers the strongest cost--benefit trade-off for many deployments. Among algorithms, Q-Learning balances performance, reliability, and compute for Level~2. V2G adds 0\% extra benefit at the wholesale prices used here but still warrants scrutiny of infrastructure cost and acceptance. Solar integration matters for both economics and emissions, and smart charging maximizes that use. Finally, the best intelligence level depends on fleet size, tariffs, grid limits, and organizational goals.

\subsection{Closing Statement}

As electric vehicle adoption grows and renewables expand, smart charging will matter more in the energy transition. Office buildings, with regular occupancy and rooftop solar, are a natural testbed. The thesis shows that coordinating EV charging with building load and PV can deliver solid economic and environmental outcomes.

Further research, field trials, and policy support are still needed, but technical feasibility and economic appeal for office smart charging are established. The open questions are less about \textit{whether} such systems can work than about \textit{when} and \textit{how} they will scale.
